\section{Lean definitions, checked results, and the iPRA export}\label{sec:formalization}

\subsection{The checked standalone kernel}

The accompanying project was compiled with Lean 4.32.0. It imports only
the distribution's \texttt{Init} and \texttt{Std} libraries; no Mathlib
download or repository checkout is required to replay it. There are
58 named proved theorems and two additional constructors whose numerical
contract fields are proved. The axiom audit prints 60 reports. Its actual
axiom union is \texttt{propext}, \texttt{Classical.choice}, and
\texttt{Quot.sound}. There are no admitted proofs or added correctness
axioms in the checked project. These facts are recorded in
\texttt{validation/lean\_status.json} and the full replay log.

\begin{longtable}{@{}p{0.25\textwidth}p{0.08\textwidth}p{0.60\textwidth}@{}}
\toprule
Module & Proofs & Exact checked scope\\
\midrule
\endhead
\texttt{FiniteMeasure} & 24 & Rational sums with internally variable finite lengths; finite Fubini; positivity; weighted errors; simultaneous perturbation of sample values and probability masses.\\
\texttt{Truncation} & 16 & Clipping bounds and nonexpansiveness, excess-tail formula and stability, scalar and finite weighted Vitali inequalities, and a three-term precision budget.\\
\texttt{Presentations} & 10 & Array-based finite probability data, refinement pullbacks, finite $L^1$ laws, and invariance of integrals, $L^1$ distance, capped distance, and tail integrals.\\
\texttt{ActionCompletion} & 8 & A concrete gain-dependent sampling schedule, regularity of mapped names, preservation of the same-name condition, and the proved fields of action composition.\\
\texttt{VitaliBlueprint} & 0 & Typechecked numerical definitions and target propositions for the stronger limit-producing theorem; the target propositions are not asserted as theorems.\\
\bottomrule
\end{longtable}

The checked namespace is \texttt{IPRAContinuousMeasure}. Representative
entry points are:
\begin{itemize}
\item expectation errors: \path{expectation_error}, \path{mass_perturbation};
\item clipping estimates: \path{finite_vitali_unit}, \path{vitali_budget};
\item refinement: \path{refinement_integral}, \path{refinement_l1},\\
      \path{refinement_rho}, \path{refinement_tail};
\item name completion: \path{action_regular}, \path{mapName_respects}.
\end{itemize}
The definitions \texttt{mapName} and
\texttt{composeDenseAction} contain checked proofs of their required
stability fields. The latter result is composition of metric actions;
linearity, positivity, the full stochastic tensor construction, and the
analytic theorems remain the report's mathematical results and development
targets rather than checked claims of this standalone core.

The project also kernel-checks four exact examples, including the
counterexample to \eqref{eq:false-old-vitali}. The separate Python programs
exercise larger finite cases. Their assertions are implementation checks;
the Lean universal statements and the mathematical proofs have separate
roles.

\subsection{All presentation data are explicit}

The interfaces are actual numerical structures rather than opaque
parameters named ``correctness.'' The following summary identifies their
fields and mathematical definitions; the full Lean definitions are included.

\begin{description}[style=nextline,leftmargin=0pt,labelwidth=0pt,labelsep=0pt]
\item[\texttt{FiniteProbability}]
An \texttt{Array Rat} of weights, a proof of nonnegativity of its in-range
entries, and a proof that its explicit rational fold sums to one. Array
access uses a total default outside the stored length. The functions
\texttt{integral}, \texttt{l1}, and \texttt{rho} are the finite sums in
\eqref{eq:m-simple} and \eqref{eq:m-stable-observers}.

\item[\texttt{Refinement}]
An \texttt{Array Nat} giving a parent for every fine atom, its range
proof, and the equation that the sum of each parent's descendant masses
equals the coarse mass. \texttt{Refinement.pull} is literal composition
of a value table with that parent array. Zero masses are allowed.

\item[\texttt{ProbabilityTower} and \texttt{SimpleTable}]
A numerical family of finite probabilities with coherent refinements;
a simple table is an array whose length matches the probability at its
stored stage. \texttt{SimpleTable.distance} lifts two arrays to the
maximum of their levels and evaluates the weighted absolute difference.
\texttt{measureDistance} and \texttt{tailIntegral} are similarly explicit.

\item[\texttt{RegularL1Name} and \texttt{SameL1Name}]
A sampler into simple tables and the bound
$d_1(s_k,s_\ell)\le2^{-k}+2^{-\ell}$.
The same-name predicate uses the corresponding cross-comparison at
arbitrary precisions. A full quotient/setoid API can be developed from
these data; it has not been silently imported as a new arithmetic sort.

\item[\texttt{DenseAction}]
A routine returning a simple output for a simple input and a precision,
a nonnegative rational gain, and the finite stability inequality.
An integer majorant for the gain is supplied separately to the name
construction and composition routines. \texttt{actionSample} uses a
precision offset based on that majorant.
The supplied implementation chooses a coarse integer offset; a proved
binary-length offset can improve its cost without changing the interface.

\item[\texttt{VitaliData}]
The uniformly named family, \texttt{UniformTailControl}, and
\texttt{InMeasureCauchyRate}, including the monotonicity and numerical
error contracts. \texttt{ceilingLog2} is a bounded search, not a real
logarithm. \texttt{vitaliPrecision}, \texttt{vitaliIndex}, and
\texttt{vitaliSample} are the actual programs in
\eqref{eq:m-vitali-rate}--\eqref{eq:m-output-name}.
\end{description}

\texttt{VitaliNameRegularityGoal}, \texttt{VitaliLiftingGoal}, and
\texttt{VitaliInheritedTailGoal} state the remaining propositions for the
constructed samples. A definition of type \texttt{Prop} records a target;
it supplies no proof of it. The analytic theorem does not assume an
integrable limit, and neither does \texttt{VitaliLiftingGoal}.

The host Lean structures permit arbitrary Lean functions as fields.
An iPRA instance fixes PR implementations for those fields and supplies
their object-theory contracts. Keeping this step explicit prevents a
higher-type semantic interface from being misreported as a uniform
first-order theorem about all arbitrary functions.

\subsection{An explicit object-language Vitali endpoint}

Here is a direct denominator-cleared target, sufficient to remove any
ambiguity about the proposed arithmetic translation. Let $w_i$ be natural
weights. Represent real samples $x_i,y_i$ with one positive denominator
$D$ and signed integer numerators $X_i,Y_i$. Let the integer cutoff
$\kappa\ge D$ represent $K=\kappa/D\ge1$. Define natural quantities
\[
 d_i=|X_i-Y_i|,\quad
 u_i=(|X_i|-\kappa)^+,\quad v_i=(|Y_i|-\kappa)^+,\quad
 c_i=\min(D,d_i).
\]
All can be written using natural truncated subtraction. If
$X_i=X_i^+-X_i^-$ and $Y_i=Y_i^+-Y_i^-$, for example,
\[
 d_i=((X_i^++Y_i^-)\mathbin{\dot-}(X_i^-+Y_i^+))
      +((X_i^-+Y_i^+)\mathbin{\dot-}(X_i^++Y_i^-)),
\]
where $a\mathbin{\dot-}b=\max(a-b,0)$; the parentheses indicate the two
complete minuends. Define the other absolute values in the same way.
The finite theorem to export is
\begin{equation}\label{eq:object-vitali}
 D\sum_{i<n}w_id_i
 \le D\sum_{i<n}w_i(u_i+v_i)
       +2\kappa\sum_{i<n}w_ic_i.
\end{equation}
Dividing by $D^2\sum_iw_i$ when that total is positive is precisely the
normalized clipping inequality. The unnormalized integer inequality
itself also holds when every weight is zero.

\paragraph{Definitions and proof program.}
Encode the four signed-numerator arrays, the weights, $n,D,\kappa$ in a
finite tuple. Define each lookup with a length guard and a zero default.
Define the three sums by primitive recursion on their prefix length $j$.
The invariant is the displayed inequality with $n$ replaced by $j$.
It is quantifier-free. At $j=0$ it is $0\le0$. At $j+1$ add the scalar
inequality
\[
 Dw_jd_j\le Dw_j(u_j+v_j)+2\kappa w_jc_j,
\]
which follows from the same finite clipping cases as
Lemma~\ref{lem:m-clipping}. Every branch compares natural integers.
The code-validity and cutoff guards remain explicit assumptions.
The theorem is then universally closed over these numerical parameters.

\paragraph{Reuse in the existing backend.}
The repository supplies the needed target syntax and rule
\texttt{PRA.qfInd} in \texttt{ConciseIPRA.PR}; natural sum/order and
signed-integer arithmetic in \texttt{IPRADerivations.Core}; rational
normalization, finite-sum semantics and monotonicity in
\texttt{Analysis.RationalFiniteSum}; and a closely related prefix
truncation proof in \texttt{Probability.TruncationEstimate}
\cite{repo-pr,repo-rational-sum,repo-truncation}.
The new file should construct a proof value of the form
\[
 [\,]\ \vdash_{\mathtt{PRAAx}}\
 \forall\vec a\,\bigl(\operatorname{Valid}(\vec a)=1
       \longrightarrow\operatorname{VitaliClear}(\vec a)=1\bigr),
\]
where \texttt{VitaliClear} is the characteristic function of
\eqref{eq:object-vitali}. This is a real proof-construction task. The
standalone theorem over \texttt{Rat} is a checked mathematical guide to it,
not an automatic quotation into \texttt{Prf PRAAx}.

\paragraph{From the finite endpoint to the named theorem.}
For fixed PR routines $S,U,R$, the desired object sentence universally
quantifies numerical indices and precisions. Its premises are the actual
regularity, tail, and capped-distance contracts of those routines. Apply
\eqref{eq:object-vitali} to finite common-refinement samples; use
\eqref{eq:m-observer-lipschitz} to account for their errors; substitute the
bounded logarithm program; and prove regularity of
$b_k=S(N(k+2),k+2)$ by the explicit three-error calculation.
The final convergence statement is expressed using cross-distances of
the two samplers with arbitrary observation errors. It need not quantify
over an uncountable space of limits. The inherited-tail property follows
by an arbitrary-precision comparison, not by a full induction schema for
analytic formulas.

\subsection{Development order for the remaining mathematics}

\begin{longtable}{@{}p{0.22\textwidth}p{0.39\textwidth}p{0.32\textwidth}@{}}
\toprule
Target & Necessary definitions and induction data & Route to completion\\
\midrule
\endhead
Full Vitali lifting & Common-refinement pseudometric laws; regular names; tail and capped-distance stability; bounded ceiling logarithm & Prove the typed goals in \texttt{VitaliBlueprint}; then implement \eqref{eq:object-vitali} and its clock substitutions in the object backend.\\
Conditional expectation & Cell numerators, exact zero-weight branch, parent-map composition, integral-characterization equation & Reuse \texttt{TowerCommonDen.regroup}; prove contraction after canceling cell mass, then lift and prove \eqref{eq:m-upward-rate}.\\
Stochastic operator category & Linear/positive/unital/integral action fields; strong-action sampler; atom-product tables & Extend the checked metric-action constructors with the stated algebraic fields. Prove equality on dense inputs first; tensor by the indicator estimate.\\
Integrated formulas & Many-sorted finite syntax with bound-variable indices; interval and Lipschitz annotations; net and quadrature codes & Structural recursion for the evaluator and budgets; induction on the fixed syntax; array induction for each numerical binder. Clip intermediate values to their annotated ranges.\\
Holomorphic integrals & Complex rational pairs; coefficient-name family; radius/majorant data; finite derivative sums & Prove finite geometric and polynomial remainder identities, then use their explicit tail clocks and the integral contraction.\\
Homological transfer & Graded matrix records, normalized contraction equations, perturbation equation, row-sum norm, finite Neumann sums & Follow the exact identity and residual route in Section~\ref{sec:homological}; use its matrix blueprint and denominator-cleared bounded folds.\\
\bottomrule
\end{longtable}

For an internally variable syntax size or matrix dimension, the compiler
must use numerical arrays and bounded characteristic folds, rather than
generate one externally unrolled theorem per size. The repository's current
bounded-syntax and rational-array providers are suitable starting points.
For a fixed mathematical formula, a schema indexed externally by that
formula is already meaningful; it should simply be labeled as such.

\subsection{Reproduction and verification limits}

Run \texttt{python3 verify.py} for the exact finite demonstrations. With
the pinned Lean toolchain installed, run \texttt{bash lean/replay.sh} from
the appropriate working directory, or use the supplied Makefile target.
\texttt{python3 build\_pdf.py} rebuilds this document. The TeX helper
can create local formats and font maps from installed TeX Live resources
on a minimally configured system; it does not download missing packages.

The manifest records the source commit and Git blob identifiers of the
upstream files actually used. Those files were inspected, not rebuilt in
this session. The standalone Lean replay, exact Python checks, and document
build were performed here. External toolchain binaries and build caches
are excluded from the archive. No claim of a complete iPRA formalization
of the whole analytic framework follows from the 58 host theorems.
